\documentclass[11pt]{article}
\usepackage[a4paper]{geometry}

% Default formatting packages
\usepackage{booktabs}
\usepackage{graphicx}
\usepackage{subcaption}
\usepackage{caption} 
\captionsetup{font=small,labelfont={bf},width=0.95\textwidth}

% Extra formatting packages
\usepackage{xcolor}
\usepackage{amssymb}
\usepackage{hyperref}
\usepackage{xstring}
\usepackage{xspace}

% Internal references, this MUST be the last package loaded, except commands below.
\usepackage{cleveref}
\usepackage[edit]{commands}

\frenchspacing

\title{The \texttt{commands} Package}
\author{Jochen Wuttke}
\date{\today}

\begin{document}

\maketitle

\begin{abstract}
The \texttt{commands} package provides custom macros for document editing and acronym handling. 
It includes debugging features for tracking first occurrences of references (in \texttt{edit} mode) 
and a smart acronym tracking system that automatically expands acronyms on first use.
\end{abstract}

\tableofcontents

\section{Introduction}

The \texttt{commands} package extends \LaTeX{} with utilities designed to improve document editing workflows:
\begin{itemize}
  \item Smart reference highlighting for identifying first uses of figures, tables, and equations
  \item Automatic acronym expansion on first occurrence
  \item Clean production output without debugging artifacts
\end{itemize}

\section{Package Options}

The package accepts the following options:

\subsection{Default Mode (No Options)}

\begin{verbatim}
\usepackage{commands}
\end{verbatim}

In default mode, only the core acronym tracking features are active. 
This is recommended for final documents and publication.

\subsection{Edit Mode}

\begin{verbatim}
\usepackage[edit]{commands}
\end{verbatim}

Edit mode enables debugging macros that highlight the first occurrence of cross-references. 
This is useful during manuscript preparation and review.

\section{Features}

\subsection{Reference Highlighting (Edit Mode Only)}

When the package is loaded with the \verb|edit| option, 
references made with referencing commands from the \verb|cleveref| package 
(\verb|\cref{}|, \verb|\Cref{}|, \verb|\crefrange{}{}|, 
\verb|\Crefrange{}{}|) are enhanced to highlight first occurrences.

\subsubsection{Behavior}

On the first use of a reference label within the document:
\begin{itemize}
  \item The reference is highlighted with a yellow background box.
  \item Subsequent uses of the same label display normally (no highlighting).
  \item This applies to both lowercase (\verb|\cref{}|) and capitalized (\verb|\Cref{}|) variants.
\end{itemize}

\subsubsection{Purpose}

Reference highlighting helps editors and authors:
\begin{itemize}
  \item Verify that all figures, tables, and equations are referenced at least once.
  \item Verify that references to floats come before the corresponding float in the print layout.
  \item Identify sections that might be missing cross-references.
\end{itemize}

\subsubsection{Example}

In edit mode, the first reference to a figure appears highlighted:

\textit{First use:} \colorbox{yellow}{Figure 1}  (appears in yellow box)

\textit{Later uses:} Figure 1 (appears normally)

\subsubsection{For Publication}

Remove the \verb|edit| option before final submission or publication:

\begin{verbatim}
\usepackage{commands}  % Remove [edit] here
\end{verbatim}

All yellow highlighting automatically disappears, leaving clean references.

\subsubsection{Range References}

The package also overrides \verb|\crefrange{}{}| and \verb|\Crefrange{}{}| 
for highlighting ranges of references (e.g., ``Figures 1--5'').

\paragraph{Syntax:}
\begin{verbatim}
\crefrange{fig:start}{fig:end}
\Crefrange{fig:start}{fig:end}
\end{verbatim}

\paragraph{Behavior:}
\begin{itemize}
  \item The range itself is highlighted on first use (e.g., the entire ``Figures 1--5'' text).
  \item Subsequent uses of the same range show normally.
  \item Both start and end labels are marked as individually referenced.
\end{itemize}

\paragraph{Example:}
\begin{verbatim}
See \crefrange{fig:start}{fig:end} for analysis.
% First use: "Figures 1--5" appears highlighted in yellow
% Second use: appears normally

Later, \cref{fig:start} and \cref{fig:end} will not be highlighted because
they were marked as referenced by the \crefrange command.
\end{verbatim}

\subparagraph*{Important Warning:}
Range references only track the explicitly mentioned start and end labels. 
Any figures that numerically fall within the range but are not explicitly referenced 
(e.g., Figure 3 in \verb|\crefrange{fig:1}{fig:5}|) 
will not be marked as referenced and will still appear highlighted when referenced individually.

However, the start and end labels will be marked as already used, and \emph{will not be highlighted again}
when used as individual labels.

This is intentional: it allows you to reference a range without making claims about all intermediate figures.

\subsection{Acronym tracking and expansion}

The \verb|\acronym| command provides intelligent acronym expansion that works in both default and edit modes.

\subsubsection{Usage}

\begin{verbatim}
\acronym{SHORTFORM}{Long Form Description}
\end{verbatim}

\paragraph{Parameters:}
\begin{description}
  \item[SHORTFORM] The acronym (e.g., \verb|XRF|, \verb|ICPMS|)
  \item[Long Form Description] The full text to display on first use 
    (e.g., \verb|X-ray fluorescence|)
\end{description}

\subsubsection{Behavior}

\paragraph{First occurrence in the document:}

The acronym command expands to:

\begin{center}
  \fbox{\textit{Long Form Description (SHORTFORM)}}
\end{center}

\paragraph{Subsequent occurrences:}

The acronym command expands to:

\begin{center}
  \fbox{\textit{SHORTFORM}}
\end{center}

This automatic tracking is done once per unique acronym in the document.

\subsubsection{Recommended Pattern}

Define a \verb|\newcommand| wrapper for each acronym you use frequently:

\begin{verbatim}
\newcommand{\XRF}{\acronym{XRF}{X-ray fluorescence}\xspace}
\newcommand{\ICPMS}{\acronym{ICPMS}{inductively coupled 
  plasma mass spectrometry}\xspace}
\end{verbatim}

Then use the command in your text:

\begin{verbatim}
Analysis was performed using \XRF and \ICPMS techniques.
% First occurrence: "X-ray fluorescence (XRF) and inductively..."
% Second occurrence: "XRF and ICPMS techniques"
\end{verbatim}

\paragraph{Note:} The \verb|\xspace| command after \verb|\acronym| provides intelligent spacing 
after the acronym, automatically handling punctuation.

\subsubsection{Example Output}

Here is how the recommended pattern works in practice:

\begin{itemize}
  \item \textbf{First paragraph:} Analysis was performed using \acronym{XRF}{X-ray fluorescence}\xspace{} 
    and \acronym{ICPMS}{inductively coupled plasma mass spectrometry}\xspace{} techniques.
  
  \item \textbf{Second paragraph:} The \acronym{XRF}{X-ray fluorescence}\xspace{} results were compared 
    with \acronym{ICPMS}{inductively coupled plasma mass spectrometry}\xspace{} measurements.
\end{itemize}

Notice that the second paragraph shows only the short forms, as the acronyms have already been introduced.

\section{Complete Example}

Below is a minimal working example demonstrating both features:

\begin{verbatim}
\documentclass{article}
\usepackage{xstring}
\usepackage{xspace}
\usepackage{cleveref}

% For editing:
\usepackage[edit]{commands}
% For publication:
% \usepackage{commands}

\newcommand{\XRF}{\acronym{XRF}{X-ray fluorescence}\xspace}
\newcommand{\TOC}{\acronym{TOC}{total organic carbon}\xspace}

\begin{document}

\section{Introduction}

We analyzed samples using \XRF and \TOC measurements.
See \cref{fig:example} for results.

\section{Results}

The \XRF data showed high concentrations.
Additionally, \TOC values indicate good preservation.
Refer to \cref{fig:example} again.

\begin{figure}
  \centering
  \caption{Sample analysis results}
  \label{fig:example}
\end{figure}

\end{document}
\end{verbatim}

In edit mode, the first \verb|\cref{fig:example}| is highlighted in yellow and uses of \verb|\XRF| and \verb|\TOC| are 
expanded to the full text description. The second uses appear normally.

\section{Implementation Notes}

\subsection{Dependencies}

The package requires:
\begin{itemize}
  \item \verb|cleveref|: For enhanced cross-references (required for edit mode)
  \item \verb|xstring|: For the \verb|\IfSubStr| command used in acronym tracking
  \item \verb|xspace|: Recommended for clean spacing after acronyms (not required)
\end{itemize}

\subsection{Load Order}

Always load packages in this order:
\begin{enumerate}
  \item Content packages (\verb|graphicx|, \verb|caption|, etc.)
  \item \verb|cleveref| (last among content packages)
  \item \verb|commands| (after \verb|cleveref|)
\end{enumerate}

\textbf{Important:} \verb|cleveref| must load before \verb|commands|, as the editing 
macros redefine its commands.

\subsection{Macro Expansion}

The acronym tracking uses \LaTeX{}'s \verb|\expandafter| and \verb|\gdef| to store 
the state of each acronym. This ensures tracking persists across multiple uses, 
even if the same macro is redefined.

\section{Troubleshooting}

\subsection{Highlighting Not Appearing}

\paragraph{Problem:} Yellow highlighting is not visible in PDF output.

\paragraph{Solution:} Ensure the package is loaded with the \verb|edit| option:
\begin{verbatim}
\usepackage[edit]{commands}  % Not: \usepackage{commands}
\end{verbatim}

\subsection{Acronym Not Tracking}

\paragraph{Problem:} The acronym expands the full form every time it is used.

\paragraph{Solution:} Verify that each acronym is unique. The tracking uses the 
acronym code (SHORTFORM) as the identifier. If you use the same code with different 
full forms, tracking will be inconsistent.

\paragraph{Also check:} Ensure you are using \verb|\acronym| directly or via a 
\verb|\newcommand| wrapper consistently. Do not mix direct and wrapper usage.

\subsection{Compilation Errors}

\paragraph{Missing Package:} If you get errors about undefined commands, ensure:
\begin{itemize}
  \item \verb|cleveref| is loaded before \verb|commands|
  \item \verb|xstring| is available (comes standard with modern \LaTeX{} distributions)
\end{itemize}

\section{Summary}

\begin{table}
  \centering
  \caption{Features by mode}
  \begin{tabular}{lcc}
    \toprule
    \textbf{Feature} & \textbf{Default} & \textbf{Edit} \\
    \midrule
    Acronym tracking & \checkmark & \checkmark \\
    Reference highlighting (\texttt{\textbackslash cref}, \texttt{\textbackslash Cref}) & & \checkmark \\
    Range reference highlighting (\texttt{\textbackslash crefrange}, \texttt{\textbackslash Crefrange}) & & \checkmark \\
    Production-ready output & \checkmark & \\
    \bottomrule
  \end{tabular}
\end{table}

Use the \verb|edit| mode during writing and review, then remove the option for final publication.

\end{document}
